\documentclass{ximera}

\ifdefined\HCode
  \newcommand{\alttext}[1]{\HCode{<span class="sr-only">#1</span>}}
\else
  \newcommand{\alttext}[1]{} % Fallback for PDF view
\fi


% MATH -----------------------------------------------------------
\newcommand{\nats}{\mathbb N}
\newcommand{\ints}{\mathbb Z}
\newcommand{\rats}{\mathbb Q}
\newcommand{\reals}{\mathbb R}
\newcommand{\complex}{\mathbb C}
\newcommand{\powerset}{\mathscr P}

%-------------------------------------------------------------

\pgfplotsset{compat=1.18}
\begin{document}
\begin{abstract}
 \end{abstract}

    \title{Exemplar Examples 24}
    
    \maketitle

Let $k\neq 0$ be constant.  Let's solve the IVP $y^{\,\prime\prime}+25y=50+k\delta \left(t-\frac{\pi}{5}\right)$, $y(0)=0$, $y^{\,\prime}(0)=1$ in terms of $k$, and express the solution as a piecewise defined function.

\begin{enumerate}
\item First let's take the Laplace transform of both sides, using $Y(s)$ to denote $\mathcal{L}{[y(t)]}$.

\vfill


\item  Let's solve for $Y(s)$.

\vfill

\item Let's re-express $Y(s)$ in a form conducive for taking the inverse Laplace transform.


\vfill

\newpage

\item  Let's find $y(t)$ by taking the inverse Laplace transform of $Y(s)$, and express the solution as a piecewise defined function.

\vfill
\vfill
\vfill

\item 

\vfill



\end{enumerate}

% Should get 2-2cos(5t)+0.2sin(5t)+0.2k U(t-pi/5)* sin[5(t-pi/5)]

% Should get 2-2cos(5t)+0.2sin(5t) - 0.2k U(t-pi/5)* sin[5t]

\end{document}